\section{Discussion}

\subsection{What the records teach about strategic behavior}

The observed agents sometimes treat a transaction as a change in competitors' future action sets. Sol prices Marvin Gardens through construction funds and possible intermediary resale. Grok refuses its own completion while discussing the proposer's ability to build before rolling. Gemini asks others to withhold cash that would fund red houses. Grok and Opus negotiate over whether another player would inherit a distressed estate. These are different expressions of one substantive pattern: the competitive consequences named in bargaining extend beyond the nominal item and its current owner.

This pattern is not dictated action by action by the Monopoly rules. The rules permit refusal, markup, accommodation, and silence; they do not require the particular resale account, cash-denial appeal, or rescue pitch. Nor does permission to bluff prescribe the \$600 final-price claim or the omitted mortgage. The prompt does shape which conduct is licensed. We therefore describe the behavior as authorized strategic interaction, without presenting deception as an emergent violation of an honesty instruction.

The most useful negative result concerns the strength of the valuation interpretation. Strategic attention has observable examples, but neither cash repricing nor pooled asking-price medians establish strategic valuation. The unchanged-bundle comparison is almost as frequent without completion, and same-deed comparisons are mixed. A mechanical haggling routine can coexist with a tailored explanation. Conversely, the generic pattern does not show that every tailored explanation is empty. The correct conclusion preserves both scales: a broad bargaining practice and local evidence of state-specific competitive reasoning. Distinguishing them is more informative than selecting the favorable price cut and naming it a premium.

Another distinction concerns public capacity versus hidden mental states. Much of the observed opponent attention concerns cash, holdings, collateral, and turn order, all visible facts. Sol revises G15 terms after an explicit objection, rather than predicting an unobserved preference and testing its calibration. G13's buyer discusses the estate but does not affirm the seller's asserted cash minimum. These interactions demonstrate attention to public competitive constraints. They do not establish theory of mind about hidden beliefs or beliefs changed by persuasion.

\subsection{Strategic effectiveness needs its own criterion}

The economic data show why the cited constraints are consequential, while also showing why later victory is an inadequate score for each decision. A completion must pass through financing and valid construction to raise its rents. Cash receipts depend on actual opponent exposure, mortgages, and the payer's solvency. Inheritance can make an earlier sale temporary, as North Carolina returns to Sol, or move a blocked group to a third player, as in G2. Evaluating a trade by its eventual owner's winning status would attribute all intervening chance and opponent action to the original bargain.

The all-debt denominator also changes the balance of evidence. Most menu-reaching debts with demonstrated payment capacity are paid. A small set of payable eliminations is meaningful, but it is not representative of every financing decision. The successful \$1,400 payment and the correct \$2,000 insolvency account deserve the same explanatory attention as the contradictions. Local debt feasibility has a defined endpoint: paying an existing obligation using permitted own assets. Bargaining quality instead depends on future responses and risks. A deterministic replay verifies a recorded history and can support a local payment sequence; it does not supply an optimal continuation policy.

First-estate concentration has a similar boundary. The first creditor often already has the largest prior rent receipts, and twelve creditors lead or tie in buildings. That supports a pathway in which productive positions precede estates. It neither makes every estate irrelevant nor identifies estate acquisition as the cause of winning. The future estate is strategically relevant to G13 and G18 because ownership can change through failure, not because the observed concentration proves a general snowball mechanism independent of prior strength.

The evidence would contradict a stronger central claim if rival-sensitive examples disappeared under exact state checks, if recipients had not received the cited statements, or if the proposed rival consequences were unrelated to the actual exchange. The corpus contains both confirming and limiting records. G2's development calculation is partly wrong, G19's recipient explicitly rejects a false premise, and G20's appeal does not attract uniform participation. Those qualifications belong to the substantive account of what the agents accomplish. They are not a separate catalog of mistakes.

\section{Relation to prior work}

Early neural negotiation work studies language and agreement under explicit item utilities~\cite{lewis2017deal}. He et al.~\cite{he2018decoupling} separate coarse strategic acts from language generation. That distinction informs our use of structured offers alongside public messages: their relationship is an empirical object, not an assumed identity. NegotiationArena supplies controlled tactic comparisons across bargaining scenarios~\cite{bianchi2024negotiationarena}. Our selected public bluffs have no comparable intervention, so we do not transfer its causal tactic interpretations to these exchanges.

Cosentino et al.~\cite{cosentino2026counterparty} directly distinguish counterparty modeling from strategic use in controlled multi-attribute bargaining, measuring outcomes against known utilities. Our distinction between attention and effectiveness is consistent with that question, while the multiplayer setting introduces construction finance and future creditor ownership. We lack its explicit utility reference. Classical work on sales with externalities formalizes how the identity of an acquirer can affect non-acquirers~\cite{jehiel1996externalities}. That provides an economic interpretation for blockers and resale concerns; we do not claim that the observed agents implement its optimal mechanism.

AgentMonopoly already studies Monopoly actions, language, model families, opponents, and persona instructions~\cite{wu2026agentmonopoly}. FinBoardBench already studies financial board games, negotiation, and liquidity, including Monopoly counteroffers~\cite{hu2026finboardbench}. Our contribution is the bounded behavioral account enabled by complete offer histories, financing opportunities, and estate trajectories, with more comparable existing-record cuts. Monopoly play, counteroffers, and language-action contrasts themselves are not new contributions.

The Vending-Bench Arena communication study joins emails to simulator state, analyzes follow-through, and repeats judging with other families~\cite{li2026arena}. It provides a useful precedent for contextual state-grounded analysis. Our eight incidents are a purposive set rather than a replicated population communication classifier. We retain the distinction between factual assertion, delivered exposure, explicit correction, and later action, while applying a different conduct norm that permits public bluffing. None of these related results validates our semantic labels or confers novelty on source bookkeeping alone.

\section{Limitations}

The twenty games have five seeds, a single cyclic opponent order, one fixed roster, visible names, and provider-default sampling. Response conditions differ by configuration and affect the entire table. Removing Opus or conflicting Grok records would not recover the unaffected game that would otherwise have occurred. We do not rank models, infer equal compute from effort settings, or infer missing reasoning from zero reported tokens.

Trade states, prices, and continuations are selected by the players. Movement exposure, ownership duration, and first-creditor status are endogenous. Board prices are scales, not utilities. The debt analysis covers episodes reaching a liquidation menu and uses inherited local feasibility results; it is not a complete policy-value analysis. No new control condition tests corrected prompts, unfamiliar opponents, or altered rules.

Semantic codes and incident selection are retrospective and LLM-assisted. Six incidents foreground Grok, whose visible response channels create unusual observability. Two same-family readers with partial unblinding do not supply independent human validation. Explanations are not privileged cognition measurements. A final publication requires accountable human authorship; this independent draft invents no human byline. These limits bound the claims rather than implying that unmeasured capability is absent.

\section{Conclusion}

The strongest supported discovery is competitive-option attention: LLM agents sometimes bargain over the finance, development, resale, and inheritance opportunities their transactions create for rivals. That attention appears in refusals of own completion, revised terms, a funding-denial appeal, and negotiations over distressed estates. It coexists with generic repricing, mixed same-deed asks, arithmetic errors, and conflicting execution. Following the surrounding economy makes the stakes visible without declaring the strategy effective. The resulting account is a study of situated strategic conduct, with optimal valuation, causal persuasion, learning, and model rankings left open.
